\subsection{Expression Vocabulary and Representation}
\label{sec:generator-vocabulary}

\subsubsection{Operators, Parameters, and Lagged Observations}

The benchmark adopts a fixed vocabulary for writing recurrence mechanisms.
Its terminals are numerical parameters and past observations; addition,
multiplication, and the hyperbolic tangent combine them into expressions.
The symbol $t$ labels the observation being generated. Each update computes
its value from the available lag history and the mechanism's fixed parameters.
Table~\ref{tab:generator-vocabulary} defines the five elements of this language.

\begin{table}[H]
\centering\small
\begin{tabular}{@{}
>{\raggedright\arraybackslash}p{0.20\linewidth}
>{\raggedright\arraybackslash}p{0.25\linewidth}
>{\raggedright\arraybackslash}p{0.46\linewidth}@{}}
\toprule
Element & Representation & Meaning \\
\midrule
Parameter & $P_a$ &
A real coefficient supplied separately, such as $P_0=0.7$. \\
Past observation & $\operatorname{LAG}(j,L_a)$ &
Reads $x_{j,t-\ell}$ after assigning the positive integer $\ell$ to $L_a$. \\
Addition & $\operatorname{ADD}(e_1,e_2)$ &
Combines two expressions as $e_1+e_2$. \\
Multiplication & $\operatorname{MUL}(e_1,e_2)$ &
Forms $e_1e_2$: coefficients, interactions, and repeated factors. \\
Hyperbolic tangent & $\operatorname{TANH}(e)$ &
Nonlinear; output in $(-1,1)$ and derivative magnitude $\leq1$. \\
\bottomrule
\end{tabular}
\caption{Selected expression vocabulary. Parameters and lagged observations
are leaves, addition and multiplication have two children, and $\tanh$ has
one child. The source index $j$ selects the observed series.}
\label{tab:generator-vocabulary}
\end{table}

\subsubsection{Abstract Syntax Trees and Parameter Binding}
An \emph{abstract syntax tree} (AST) describes operations and their
arguments; a binding supplies numerical coefficients and positive integer
lags. Structure and numerical values are retained separately. Repeated
trajectories preserve both the tree and its binding while resampling initial
conditions or innovations according to the selected policy.

Let $\mathcal E$ denote any valid expression, including a complete output
tree or one of its subtrees. The recursive grammar is
\begin{equation}
\begin{aligned}
\mathcal E ::= {}& P_a \mid \operatorname{LAG}(j,L_a)\\
 &\mid\operatorname{ADD}(\mathcal E_1,\mathcal E_2)
  \mid\operatorname{MUL}(\mathcal E_1,\mathcal E_2)\\
 &\mid\operatorname{TANH}(\mathcal E_1).
\end{aligned}
\label{eq:generator-vocabulary-grammar}
\end{equation}
Here $\mathcal E_1$ and $\mathcal E_2$ are expressions built by the same rules.
The tree stores named coefficient slots $P_a$ and lag slots $L_a$; a binding
assigns real coefficients and positive integer delays. For example,
$P_0\operatorname{LAG}(0,L_0)+P_1$ becomes
$0.7x_{0,t-2}+0.01$ when $P_0=0.7$, $L_0=2$, and $P_1=0.01$.
Negative coefficients express subtraction, and repeated lag factors express
powers such as $x_{0,t-1}^2$. One output tree is retained per target series;
its lag leaves can refer to itself or another series in the system.

\paragraph{Temporal evolution through recurrence.}
Growth can arise from repeated updates of the state. For the deterministic
example
\begin{equation}
 x_t=x_{t-1}+7,\qquad t\geq1,
 \label{eq:generator-recurrent-trend}
\end{equation}
each update adds seven to the previous value. Starting from $x_0$, the
sequence is $x_0,x_0+7,x_0+14,\ldots$. Thus a trend can be represented using
only a lag leaf, a constant, and addition. This example has lag derivative
equal to one: the strict contraction conditions in
Section~\ref{sec:generator-stability} select a narrower set of mechanisms.
The shared vocabulary describes expressions; each generation family adds
its own admissibility conditions.

\paragraph{Innovations, initialization, and method-specific subsets.}
The AST represents the deterministic part of the recurrence. The executor
adds the innovation $\varepsilon_{k,t}$ according to the selected noise law,
and the initial history supplies the observations required by the lags.
These choices are recorded alongside the expression. The original sampler
and bounded-domain constructor use sums and products. The Shao--Wu
construction (Section~\ref{sec:implemented-shao-wu}) wraps lag leaves in
$\tanh$ within nonlinear products, or applies $\tanh$ once after each
complete product, according to the selected mode. The first choice checks
global moment contraction; the second checks a uniform amplitude moment
bound. Heaps
(Section~\ref{sec:implemented-heaps}) expresses its linear VAR coefficients
using parameters, lag leaves, sums, and products between coefficients and
observations.

\paragraph{Alignment with the current implementation.}
The corresponding node classes and bindings are defined in
\path{src/tsfmxai/representations/expression.py}. The base sampler uses
parameters, lags, addition, and multiplication. Specialized constructors can
also introduce $\tanh$, as shown in Table~\ref{tab:generator-vocabulary}.

\paragraph{Size of the encoder input.}
The Transformer encoder that maps an AST to its canonical form reads the tree
as a prefix sequence in which every node (a binary operator, a coefficient, a
lag leaf, or a $\tanh$) is one token, and each position in the sequence has its
own learned position embedding. The hyperparameter \texttt{max\_nodes} is the
size of that table: it is the \emph{largest number of nodes an input tree may
have}, and a longer tree cannot be encoded. A sampled tree with $K$ binary
operations has $2K+1$ nodes. Product-lag protection may additionally wrap each
of its at most $K+1$ lag leaves in a $\tanh$, so a protected tree has at most
$3K+2$ nodes. With the default limit of $K=6$ operations this gives 13 nodes
for a raw tree and at most 20 for a protected one, hence
\texttt{max\_nodes}${}=20$ in
\path{configs/canonicalization_5m_raw_tanh.json}. The bound does not depend on
the lag range, but the vocabulary size does: the codec has
$2L+5$ symbols for $L$ lags and one series (37 for $L=16$). The model is
defined in \path{src/tsfmxai/symbolic_regression/canonical_models.py}.

\paragraph{Trained encoder--decoder.}
Two encoder--decoders were trained on a corpus of $10^7$ sampled trees
(1--6 operations, lags 1--16, degree at most 4, protection probability $0.5$).
Skeletons are split by a hash of the canonical skeleton into 70/10/20 percent
for training, validation, and test; validation and test keep one row per
skeleton and no skeleton occurs in two partitions. The stored corpus has
9{,}856{,}116 rows: 6{,}997{,}384 for training, 952{,}965 for validation, and
1{,}905{,}767 for test. The \emph{skeleton} model writes only the symbols of the
canonical form; the \emph{joint} model also writes one coefficient at every
end-of-term symbol. Both use a two-layer Transformer encoder of width 96, a GRU
decoder with attention over the tree, ten epochs, and batch size 256, and were
trained on one GPU in 4.9 hours. Table~\ref{tab:canonical-training} reports free
decoding on the whole test partition; a coefficient is correct when it lies
within an absolute and relative tolerance of $0.01$.

\begin{table}[htbp]
\centering
\caption{Free decoding on the 1{,}905{,}767 test skeletons, none of which occurs
in training. ``Joint exact'' requires the skeleton and every coefficient to be
correct.}
\label{tab:canonical-training}
\begin{tabular}{lcc}
\toprule
Metric & Skeleton model & Joint model\\
\midrule
Skeleton exact & 0.9794 & 0.9732\\
Valid output & 0.9917 & 0.9869\\
Term F1 & 0.9956 & 0.9950\\
Joint exact & --- & 0.0420\\
Mean absolute coefficient error (matched terms) & --- & 0.0348\\
\bottomrule
\end{tabular}
\end{table}

Both models reached their best validation result at the last epoch and the
validation loss was still decreasing, so the ten-epoch budget, not convergence,
ended training, and the values are lower bounds for this architecture. The
canonical structure is learned well; the coefficients, which are sums of
products of several tree parameters, are not recovered precisely. The executed
record is \path{experiments/canonicalization/06_train.ipynb}, and the trained
models and the corpus are published on Hugging Face as
\path{mateuswiteck/tsfmxai-canonicalization-encoder-decoder} and
\path{mateuswiteck/tsfmxai-canonicalization-asts}.
